% File: tikz/FA22FurtherHypothesisTestsVarianceEnrichmentTikZ-001.tex
% Asset ID: FA22FurtherHypothesisTestsVarianceEnrichmentTikZ-001
% Unit code: FA22
% Topic ID: FA22FurtherHypothesisTestsVarianceEnrichment
% Related lesson file: FA22_further_hypothesis_tests_variance_enrichment_lesson.md
% Related lesson section: Section 9 Visual Asset Integration
% Used placeholder:
% [VISUAL PLACEHOLDER: FA22FurtherHypothesisTestsVarianceEnrichmentTikZ-001 | Source: Teacher transcript confidence interval derivation | Insert from tikz/FA22FurtherHypothesisTestsVarianceEnrichmentTikZ-001.tex | Purpose: Precise mathematical number-line diagram showing reciprocal inequality reversal when solving for sigma^2.]
% Source: Teacher transcript confidence interval derivation
% Source status: Off-spec enrichment, not CCEA core
% Purpose:
% Precise mathematical number-line diagram showing why reciprocal inequalities reverse
% when solving for the population variance sigma^2 in a chi-squared confidence interval.
%
% Boundary note:
% This asset supports an off-spec enrichment lesson. It must not be used to imply
% that chi-squared confidence intervals for variance are required CCEA FA22 content.

\documentclass[tikz,border=8pt]{standalone}

\usepackage{amsmath}
\usepackage{amssymb}
\usepackage[T1]{fontenc}
\usetikzlibrary{arrows.meta,positioning,calc,decorations.pathreplacing}

\definecolor{Canvas}{HTML}{FAF9F6}
\definecolor{Panel}{HTML}{FFFFF0}
\definecolor{Line}{HTML}{E5E5EA}
\definecolor{TextDark}{HTML}{2C2C2E}
\definecolor{Gold}{HTML}{C5A059}
\definecolor{SoftGold}{HTML}{D4AF37}
\definecolor{Blush}{HTML}{FBEFEF}

\begin{document}

\begin{tikzpicture}[
    font=\sffamily,
    every node/.style={text=TextDark},
    box/.style={rounded corners=10pt, draw=Line, line width=1pt, fill=Panel, inner sep=12pt, align=center},
    goldbox/.style={rounded corners=10pt, draw=Line, line width=1pt, fill=Blush, inner sep=12pt, align=center},
    arrow/.style={-{Latex[length=3mm,width=2mm]}, line width=1.1pt, draw=Gold},
    emphasis/.style={draw=Gold, line width=1.2pt, rounded corners=8pt, fill=Blush}
]

\fill[Canvas] (-8.4,-6.1) rectangle (8.4,5.7);
\draw[Line, line width=1.2pt, rounded corners=18pt, fill=Panel] (-8.0,-5.7) rectangle (8.0,5.3);

\node[align=center, font=\sffamily\bfseries\Large] at (0,4.85) {Reciprocal Inequality Reversal in the Variance Confidence Interval};
\node[align=center, font=\sffamily\small, text=TextDark] at (0,4.45) {Off-spec enrichment: the algebraic trapdoor when solving for $\sigma^2$};

\node[box, minimum width=14.0cm] (defs) at (0,3.62) {
    Let
    \[
        \nu=n-1,\qquad
        A=\chi^2_{\nu}(0.975),\qquad
        B=\chi^2_{\nu}(0.025).
    \]
    With upper-tail notation, \(A\) is the smaller critical value and \(B\) is the larger critical value.
};

\node[goldbox, minimum width=13.6cm] (row1) at (0,2.25) {
    Start from the middle probability statement:
    \[
        A
        <
        \frac{\nu S^2}{\sigma^2}
        <
        B
    \]
};

\draw[Line, line width=1pt] (-5.8,1.38) -- (5.8,1.38);
\foreach \x/\lab in {-4.2/$A=\chi^2_{\nu}(0.975)$,0/$\dfrac{\nu S^2}{\sigma^2}$,4.2/$B=\chi^2_{\nu}(0.025)$}
{
    \draw[Gold, line width=1pt] (\x,1.22) -- (\x,1.54);
    \node[font=\sffamily\small, align=center] at (\x,0.92) {\lab};
}
\draw[SoftGold, line width=2pt] (-4.2,1.38) -- (4.2,1.38);
\node[font=\sffamily\small, text=Gold] at (0,1.72) {Increasing order before reciprocals};

\draw[arrow] (0,0.62) -- (0,0.05);
\node[font=\sffamily\small, text=Gold, fill=Panel, inner sep=3pt] at (2.55,0.34) {take reciprocals};

\node[box, minimum width=13.6cm] (row2) at (0,-0.72) {
    Reciprocals reverse the inequality signs:
    \[
        \frac{1}{A}
        >
        \frac{\sigma^2}{\nu S^2}
        >
        \frac{1}{B}
    \]
};

\draw[arrow] (0,-1.45) -- (0,-2.02);
\node[font=\sffamily\small, text=Gold, fill=Panel, inner sep=3pt] at (2.52,-1.74) {rewrite in increasing order};

\node[goldbox, minimum width=13.6cm] (row3) at (0,-2.78) {
    In increasing order:
    \[
        \frac{1}{B}
        <
        \frac{\sigma^2}{\nu S^2}
        <
        \frac{1}{A}
    \]
};

\draw[Line, line width=1pt] (-5.8,-3.62) -- (5.8,-3.62);
\foreach \x/\lab in {-4.2/$\dfrac{1}{B}$,0/$\dfrac{\sigma^2}{\nu S^2}$,4.2/$\dfrac{1}{A}$}
{
    \draw[Gold, line width=1pt] (\x,-3.78) -- (\x,-3.46);
    \node[font=\sffamily\small, align=center] at (\x,-4.08) {\lab};
}
\draw[SoftGold, line width=2pt] (-4.2,-3.62) -- (4.2,-3.62);
\node[font=\sffamily\small, text=Gold] at (0,-3.26) {Increasing order after reciprocals};

\node[emphasis, minimum width=13.8cm] (final) at (0,-5.0) {
    Multiply throughout by \(\nu S^2=(n-1)S^2\):
    \[
        \boxed{
        \frac{(n-1)S^2}{\chi^2_{n-1}(0.025)}
        <
        \sigma^2
        <
        \frac{(n-1)S^2}{\chi^2_{n-1}(0.975)}
        }
    \]
};

\node[rounded corners=10pt, draw=Line, fill=TextDark, text=Panel, align=center, font=\sffamily\small, inner sep=9pt, text width=4.4cm] at (5.45,0.02) {
    Memory ember:\\
    \(0.025\) is the right-tail probability for the large \(\chi^2\) value, so it sits in the lower-bound denominator.
};

\draw[decorate, decoration={brace, amplitude=6pt, mirror}, Gold, line width=1pt] (-5.7,-5.52) -- (5.7,-5.52) node[midway, below=8pt, font=\sffamily\small, text=Gold] {confidence interval for the population variance};

\end{tikzpicture}

\end{document}
